\documentclass[11pt]{article}
\usepackage[T1]{fontenc}
\usepackage[utf8]{inputenc}
\usepackage{lmodern,microtype}
\usepackage[a4paper,top=22mm,bottom=23mm,left=23mm,right=23mm]{geometry}
\usepackage{amsmath,amssymb,amsthm,mathtools,mathrsfs}
\usepackage{booktabs,longtable,array,tabularx,multirow}
\usepackage[table]{xcolor}
\usepackage{graphicx,tikz}
\usetikzlibrary{arrows.meta,positioning,calc,fit}
\usepackage{needspace,fancyhdr,titlesec,enumitem}
\usepackage[most]{tcolorbox}
\usepackage{hyperref,xurl}
\definecolor{Navy}{HTML}{244F7B}
\definecolor{Teal}{HTML}{167C80}
\definecolor{Amber}{HTML}{B77623}
\definecolor{Ice}{HTML}{F3F6FA}
\definecolor{Ink}{HTML}{202B36}
\hypersetup{colorlinks=true,linkcolor=Navy,citecolor=Teal,urlcolor=Navy,
  pdftitle={Gaussian rMPS Moments: Two Self-Contained Proof Chains},
  pdfauthor={Method Problems -- consolidated research note},bookmarksnumbered=true}
\color{Ink}
\pagestyle{fancy}\fancyhf{}
\fancyhead[L]{\sffamily\small Gaussian rMPS moments}
\fancyhead[R]{\sffamily\small Two proof chains}
\fancyfoot[L]{\sffamily\footnotesize Consolidated research note\;\textbullet\;16 September 2026}
\fancyfoot[R]{\sffamily\small\thepage}
\renewcommand{\headrulewidth}{.4pt}
\setlength{\headheight}{15pt}
\setlength{\emergencystretch}{2em}
\setlength{\parskip}{3pt}
\setlength{\parindent}{0pt}
\linespread{1.04}
\allowdisplaybreaks[1]
\clubpenalty=10000\widowpenalty=10000
\titleformat{\section}{\Large\sffamily\bfseries\color{Navy}}{\thesection}{.65em}{}
\titleformat{\subsection}{\large\sffamily\bfseries\color{Navy}}{\thesubsection}{.6em}{}
\titleformat{\subsubsection}{\normalsize\sffamily\bfseries}{\thesubsubsection}{.6em}{}
\titlespacing*{\section}{0pt}{2.0ex plus 1ex minus .2ex}{1.0ex}
\titlespacing*{\subsection}{0pt}{1.8ex plus 1ex minus .2ex}{.7ex}
\numberwithin{equation}{section}
\newtheorem{theorem}{Theorem}[section]
\newtheorem{lemma}[theorem]{Lemma}
\newtheorem{proposition}[theorem]{Proposition}
\newtheorem{corollary}[theorem]{Corollary}
\theoremstyle{definition}\newtheorem{definition}[theorem]{Definition}
\theoremstyle{remark}\newtheorem{remark}[theorem]{Remark}
\newcommand{\R}{\mathbb R}\newcommand{\E}{\mathbb E}\newcommand{\Prob}{\mathbb P}
\newcommand{\tr}{\operatorname{tr}}\newcommand{\Ent}{\operatorname{Ent}}
\newcommand{\norm}[1]{\lVert#1\rVert}\newcommand{\ip}[2]{\langle#1,#2\rangle}
\newcommand{\one}{\mathbf1}\newcommand{\F}{\mathbb F}
\newcommand{\step}[1]{\par\addvspace{.65\baselineskip}\Needspace{4\baselineskip}\noindent\textbf{#1}\par\nobreak\smallskip}
\newcolumntype{Y}{>{\raggedright\arraybackslash}X}
\newcolumntype{P}[1]{>{\raggedright\arraybackslash}p{#1}}
\newtcolorbox{takeaway}[1][]{colback=Ice,colframe=Navy!50,boxrule=.45pt,
  arc=1pt,left=8pt,right=8pt,top=7pt,bottom=7pt,fonttitle=\sffamily\bfseries,
  coltitle=Navy,colbacktitle=Ice,enhanced,#1}
\setlist{nosep,leftmargin=1.5em}

\newcommand{\C}{\mathbb C}
\newcommand{\diag}{\operatorname{diag}}
\newcommand{\cyc}{\operatorname{cyc}}
\newcommand{\cF}{c_{\F,p}}

\begin{document}
\thispagestyle{empty}
{\sffamily\color{Navy}\small METHOD PROBLEMS\hfill PROOF-CHAIN MAP\par}
\vspace{3mm}
{\sffamily\bfseries\fontsize{24}{28}\selectfont Gaussian rMPS moments\par}
\vspace{1.5mm}
{\sffamily\large Where the proof chains split and meet\par}
\vspace{2mm}
{\small Read downward. Teal and blue show the two routes; full-width boxes\par
show the shared input, convergence, and common completion.\par}
\vspace{2mm}

% This diagram is drawn at its final physical size: no resizebox or scaling.
% All vertical placements use node anchors, so text cannot outgrow a fixed
% centre-to-centre spacing. The complete mathematical proofs follow on page 2.
\newlength{\RMPSFullWidth}
\setlength{\RMPSFullWidth}{\dimexpr\textwidth-16pt\relax}
\newlength{\RMPSColumnWidth}
\setlength{\RMPSColumnWidth}{.5\textwidth}
\addtolength{\RMPSColumnWidth}{-5mm}
\addtolength{\RMPSColumnWidth}{-14pt}
\newlength{\RMPSColumnOffset}
\setlength{\RMPSColumnOffset}{.25\textwidth}
\addtolength{\RMPSColumnOffset}{2.5mm}
\begin{center}
\begin{tikzpicture}[
  >={Latex[length=1.6mm,width=1.25mm]},
  every node/.style={font=\fontsize{10.2}{12.4}\selectfont},
  base/.style={rounded corners=1.5pt,line width=.5pt,
    align=center,inner xsep=6pt,inner ysep=5pt,
    outer sep=.5pt},
  shared/.style={base,draw=Navy!70,fill=Ice,text width=\RMPSFullWidth},
  abox/.style={base,draw=Teal!85,fill=Teal!4,
    text width=\RMPSColumnWidth,minimum height=18mm},
  bbox/.style={base,draw=Navy!75,fill=Navy!3,
    text width=\RMPSColumnWidth,minimum height=18mm},
  ahead/.style={base,draw=Teal,fill=Teal,text=white,
    text width=\RMPSColumnWidth,minimum height=10mm,
    font=\sffamily\bfseries\fontsize{10.4}{12.2}\selectfont},
  bhead/.style={ahead,draw=Navy,fill=Navy},
  scope/.style={shared,draw=Amber!80,fill=Amber!7},
  flow/.style={->,draw=Navy!80,line width=.6pt},
  aflow/.style={->,draw=Teal!90,line width=.6pt}
]
\node[shared,minimum height=17mm] (model)
  {\textbf{Same Gaussian probes; same fixed subspace}\\[2pt]
   $y_a=U^*\omega_a,\quad X_a=y_ay_a^*-I_2,
      \quad G=N^{-1}\sum_{a=1}^{N}X_a$};

\node[ahead,below=7mm of model.south,
  xshift=-\RMPSColumnOffset] (ah) {CHAIN A\\Wick / transfer extension};
\node[bhead,below=7mm of model.south,
  xshift=\RMPSColumnOffset] (bh) {CHAIN B\\Multiplication / energy};
\coordinate (split) at ($(model.south)+(0,-3mm)$);
\draw[aflow] (model.south)--(split)-|(ah.north);
\draw[flow] (model.south)--(split)-|(bh.north);

\node[abox,below=3.5mm of ah] (a1)
  {\textbf{A1. Integrate each core jointly}\\[2pt]
   Pairings of $2p$ slots (real);\\
   permutations in $S_p$ (complex)};
\node[bbox,below=3.5mm of bh] (b1)
  {\textbf{B1. Use original-law multiplication}\\[2pt]
   Express $\E\tr G^4$ as a vacuum energy;\\
   repeated cores are never resampled};
\draw[aflow] (ah.south)--(a1.north);
\draw[flow] (bh.south)--(b1.north);

\node[abox,below=4mm of a1] (a2)
  {\textbf{A2. Retain the bond transfer states}\\[2pt]
   $T_p(\alpha,\beta)=\chi^{\ell(\alpha,\beta)-p}$\\
   Use cycle counts in the complex case};
\node[bbox,below=4mm of b1] (b2)
  {\textbf{B2. Restrict to reachable states}\\[2pt]
   Local Hermite degree $\le4$ per core;\\
   at most two active probe blocks};
\draw[aflow] (a1.south)--(a2.north);
\draw[flow] (b1.south)--(b2.north);

\node[abox,below=4mm of a2] (a3)
  {\textbf{A3. Sum the transfer weights}\\[2pt]
   $T_p\one=\rho_p(\chi)\one$\\
   Each coefficient network costs $\le\|u\|^{2p}$};
\node[bbox,below=4mm of b2] (b3)
  {\textbf{B3. Contract one Gaussian core}\\[2pt]
   $\E_R\|RC\|_{\mathrm F}^{2p}\le
       \rho_p(\chi)\|C\|_{\mathrm F}^{2p}$\\
   Keep the coefficient Gram $CC^*$};
\draw[aflow] (a2.south)--(a3.north);
\draw[flow] (b2.south)--(b3.north);

\node[abox,below=4mm of a3] (a4)
  {\textbf{A4. Propagate through the bonds}\\[2pt]
   $\one^{\mathsf T}T_p^{n-1}\one=c_{\F,p}\mu_p$\\
   Product tensors attain the bound};
\node[bbox,below=4mm of b3] (b4)
  {\textbf{B4. Iterate conditional energies}\\[2pt]
   $\E S_j^p\le\rho_p(\chi)\E S_{j-1}^p$\\
   The final core contributes $c_{\F,p}$};
\draw[aflow] (a3.south)--(a4.north);
\draw[flow] (b3.south)--(b4.north);

% A fit node provides a true lower boundary even if one column becomes taller.
\node[fit=(a4)(b4),inner sep=0pt,outer sep=0pt] (lastrow) {};
\node[shared,below=7mm of lastrow.south,minimum height=24mm] (scalar)
  {\textbf{First convergence: the same sharp scalar moment bound}\\[2pt]
   $\E|u^*\omega|^{2p}\le c_{\F,p}\mu_p\|u\|^{2p}$\\[3pt]
   $\rho_p(\chi)=\prod_{j=0}^{p-1}(1+\eta j/\chi),
      \quad \mu_p=\rho_p(\chi)^{n-1}$\\[1pt]
   $\eta=2$ and $c_{\R,p}=(2p-1)!!$ (real);
   $\eta=1$ and $c_{\C,p}=p!$ (complex)};
\draw[aflow] (a4.south)--++(0,-3mm)-|([xshift=-24mm]scalar.north);
\draw[flow] (b4.south)--++(0,-3mm)-|([xshift=24mm]scalar.north);

\node[shared,below=4.5mm of scalar,minimum height=24mm] (assembly)
  {\textbf{Shared completion: projected moments and sample assembly}\\[3pt]
   $\displaystyle\E\tr G^4=
   \frac{\E\tr X^4+(N-1)\bigl[2\tr V^2+\E\tr(XX'XX')\bigr]}{N^3}$\\[3pt]
   $V=\E X^2$; $X'$ is an independent copy. Keep the alternating word.};
\draw[flow] (scalar.south)--(assembly.north);

\node[scope,below=4.5mm of assembly,minimum height=17mm] (endpoint)
  {\textbf{Same endpoint; identical fixed-order constants}\\[2pt]
   Uniform bound for every fixed two-dimensional $U$;\\
   exact formula for the product-prefix subspace};
\draw[flow] (assembly.south)--(endpoint.north);
\end{tikzpicture}
\end{center}
\clearpage

\section{Model, theorem, and comparison at a glance}\label{sec:model}
\subsection{The unchanged Gaussian rMPS model}
Fix integers $n\ge2$, $d\ge2$, $\chi\ge1$, and $N\ge1$. The tensor length is $n$,
the physical dimension is $d$, the bond dimension is $\chi$, and the number of
independent probes is $N$. The ambient dimension is $d^n$, not $N$.
Write $\F\in\{\R,\C\}$, with $z\sim N_{\C}(0,1)$ meaning
$z=(g+ih)/\sqrt2$ for independent real standard Gaussians.

Let the independent standard Gaussian cores have sizes
\[
 H_1\in\F^{d\times\chi},\qquad
 H_j\in\F^{\chi\times d\times\chi}\ (2\le j<n),\qquad
 H_n\in\F^{\chi\times d}.
\]
Define the random tensor by
\begin{equation}\label{eq:model}
 \omega(i_1,\ldots,i_n)=\chi^{-(n-1)/2}
 \sum_{\alpha_1,\ldots,\alpha_{n-1}=1}^{\chi}
 H_1(i_1,\alpha_1)
 \left[\prod_{j=2}^{n-1}H_j(\alpha_{j-1},i_j,\alpha_j)\right]
 H_n(\alpha_{n-1},i_n).
\end{equation}
The product over interior sites is empty when $n=2$.
This is the normalization of Definition~1.1 of Cama\~no--Epperly--Meyer--Tropp
\cite{CEMT}, with its deterministic scaling collected into a single factor.
Gaussian tensor-train random projections already appear in Rakhshan--Rabusseau
\cite{RR2020}; the model is not introduced in this note.
Equivalently, multiply each of the first $n-1$ standard cores by $\chi^{-1/2}$
and leave the last one unscaled. Both descriptions give the same random tensor.

Fix a deterministic isometry $U:\F^2\to\F^{d^n}$, and let $\omega_1,\ldots,\omega_N$
be independent copies of \eqref{eq:model}. Set
\begin{equation}\label{eq:G}
 y_a=U^*\omega_a,\qquad Y_a=y_ay_a^*,\qquad X_a=Y_a-I_2,
 \qquad G=\frac1N\sum_{a=1}^N X_a.
\end{equation}
Every statement is for \emph{each fixed} $U$, with constants independent of its
coefficients. It is not a simultaneous statement over all subspaces chosen after
observing the probes. Unmarked matrix norms are operator norms.

Direct contraction of two copies of the cores gives
\begin{equation}\label{eq:isotropy}
 \E\omega\omega^*=I_{d^n},\qquad \E Y_a=I_2,\qquad \E X_a=0.
\end{equation}
Indeed, the physical and bond indices of the two copies must agree. The
$\chi^{n-1}$ surviving bond paths cancel the squared normalization.

\subsection{Moment factors and the two shared endpoints}
For integers $p\ge1$, define
\begin{equation}\label{eq:rho}
 \eta=\begin{cases}2,&\F=\R,\\1,&\F=\C,\end{cases}\qquad
 \rho_p(\chi)=\prod_{j=0}^{p-1}\left(1+\frac{\eta j}{\chi}\right),\qquad
 \mu_p=\rho_p(\chi)^{n-1}.
\end{equation}
Set $\rho_0=\mu_0=1$. In particular, $\mu_1=1$.

\begin{theorem}[Uniform two-dimensional endpoint]\label{thm:uniform}
For every fixed two-dimensional isometry $U$, the real-core model satisfies
\begin{equation}\label{eq:uniformR}
 \boxed{\displaystyle
 \E\tr G^4\le
 \frac{384\mu_4+3(N-1)(8\mu_2-2)^2}{N^3}.}
\end{equation}
For circular complex Gaussian cores,
\begin{equation}\label{eq:uniformC}
 \boxed{\displaystyle
 \E\tr G^4\le
 \frac{120\mu_4+3(N-1)(6\mu_2-2)^2}{N^3}.}
\end{equation}
Both proof chains give these same constants. No coefficient-sign condition,
common Schmidt basis, or bound on the Schmidt ranks of $U$ is required.
\end{theorem}

Define the product-prefix subspace
\begin{equation}\label{eq:prefixU}
 U_0e_1=e_1^{\otimes(n-1)}\otimes e_1,\qquad
 U_0e_2=e_1^{\otimes(n-1)}\otimes e_2.
\end{equation}
Its two coordinates share the first $n-1$ cores; their dependence is not discarded.

\begin{theorem}[Exact product-prefix endpoint]\label{thm:exact}
For $U=U_0$, write $a=\mu_2$, $b=\mu_3$, and $c=\mu_4$. Then
\begin{equation}\label{eq:exact}
 \boxed{\displaystyle\E\tr G^4=\frac{A_{\F}+(N-1)B_{\F}}{N^3},}
\end{equation}
where
\begin{align}
 A_{\R}&=384c-192b+48a-6,
 &B_{\R}&=88a^2-48a+6,\label{eq:ABR}\\
 A_{\C}&=120c-96b+36a-6,
 &B_{\C}&=48a^2-36a+6.\label{eq:ABC}
\end{align}
The corresponding second-moment identities are
\begin{equation}\label{eq:second}
 \E\tr G^2=\begin{cases}(8a-2)/N,&\F=\R,\\(6a-2)/N,&\F=\C.\end{cases}
\end{equation}
\end{theorem}

\subsection{What is actually being compared}
\begingroup\small\setlength{\tabcolsep}{5pt}\renewcommand{\arraystretch}{1.24}
\begin{tabularx}{\textwidth}{@{}P{.19\textwidth}YY@{}}
\toprule
\textbf{Feature}&\textbf{Chain A: Wick / transfer}&\textbf{Chain B: multiplication / energy}\\
\midrule
Local description&Pairings or permutations for repeated Gaussian occurrences.&The original-law vacuum state, core degrees, and retained coefficient matrices.\\
Bond estimate&Exact transfer row sum $\rho_p(\chi)$.&Conditional Frobenius-energy factor $\rho_p(\chi)$.\\
Coefficient control&Contract the coefficient network and apply Frobenius Cauchy--Schwarz.&Form $CC^*$, then apply convexity to its eigenvalue weights.\\
Common output&\multicolumn{2}{P{.74\textwidth}}{Sharp scalar moments; spherical averaging; the exact ordered fourth-moment identity for independent sample blocks.}\\
Explicit constants&\multicolumn{2}{P{.74\textwidth}}{Identical in both endpoint theorems and all reported fourth-moment certificates.}\\
Limitation&The source paper's fourth-moment surrogate cannot simply be used at eighth order.&The graded description is compatible with the proof, but its specialized machinery is not necessary for this fixed-order result.\\
\bottomrule
\end{tabularx}
\endgroup

\subsection{Reading guide and attribution}
Section~\ref{sec:A} develops the higher-order Wick/transfer route;
Section~\ref{sec:B} gives the multiplication/energy route. They meet at the scalar
moment theorem. Section~\ref{sec:shared} completes the common projected-moment and
sample-label steps. Section~\ref{sec:exact} computes the exact benchmark, and
Section~\ref{sec:comparison} compares constants and numerical certificates.

The two routes have identifiable public antecedents. Rakhshan--Rabusseau
\cite[Theorem~1 and Section~5.1, Lemmas~3--4]{RR2020} analyze Gaussian tensor-train
projections using a fourth-order one-core Frobenius contraction and its iteration
along the tensor train. This is the direct predecessor of the conditional-energy
calculation in Chain~B. Its extension to every integer moment order is proved
below; the one-core strategy itself is not claimed as new.

Cama\~no--Epperly--Meyer--Tropp
\cite[Section~4.4.1, equation~(4.8), and Theorems~4.3--4.4]{CEMT}
obtain the exact fourth-moment tensor and its fractured-Gaussian surrogate;
Section~4.6 treats the complex case. Chain~A extends their original-core
contraction methodology to the higher orders needed here. That extension is
not attributed to a published eighth-order theorem in their paper.

The local Wick identities are classical \cite{Isserlis,Goodman}; the transfer
kernels are normalized orthogonal/unitary pairing Gram matrices
\cite[Section~2]{CollinsMatsumoto}. Chain~B's multiplication and energy interfaces
follow the supplied v11 manuscript \cite{Framework}, while its Gaussian
conditioning inputs and the sample-label expansion have independent classical
proofs. The contribution of this note is the explicit higher-order calculation,
the specified centered-Gram formulas, and their comparison, not origination of
these ingredients. This attribution audit identifies directly relevant
antecedents; it is not an exhaustive publication-priority determination.

\par\Needspace{12\baselineskip}
\section{Chain A: the higher-order Wick/transfer proof}\label{sec:A}
\begin{takeaway}
The real eighth moment has $105$ local pairing states; the complex eighth moment
has $24$ permutation states. Their transfer matrices have constant row sums that
can be evaluated exactly. No norm estimate of a large coefficient tensor is needed
to sum those weights.
\end{takeaway}

\subsection{A1. Integrate repeated occurrences of a core jointly}
For a real centered Gaussian family, the Wick--Isserlis rule writes an even
moment as a sum over pairings, with each pair contributing its covariance
\cite{Isserlis}. For independent circular complex Gaussians, applying that rule
to real and imaginary parts leaves only pairings of an unconjugated occurrence
with a conjugated occurrence; balanced moments are indexed by permutations.
Goodman \cite{Goodman} supplies the classical complex Gaussian/Wishart setting.
These Gaussian moment identities are not new ingredients.

Expanding $|u^*\omega|^{2p}$ produces $2p$ occurrences of each original core.
Integrating that core once, with all its occurrences present, yields the local
states
\[
 \mathcal P_2(2p)\quad\text{over }\R,\qquad
 S_p\quad\text{over }\C.
\]
Their cardinalities are $(2p-1)!!$ and $p!$, respectively.
No repeated core is replaced by a fresh independent copy.

\subsection{A2. Form the bond transfer matrix}
For real pairings $\alpha,\beta\in\mathcal P_2(2p)$, let
$\ell(\alpha,\beta)$ be the number of components in their union.
A pair shared by both matchings is counted as one component. Each component
leaves one freely summable bond index, contributing a factor $\chi$.
The normalization from the $2p$ copies contributes $\chi^{-p}$ per bond. Thus
\begin{equation}\label{eq:Treal}
 T_p(\alpha,\beta)=\chi^{\ell(\alpha,\beta)-p}.
\end{equation}
In the complex model the analogous transfer is
\begin{equation}\label{eq:Tcomplex}
 T_p^{\C}(\sigma,\tau)
 =\chi^{\cyc(\sigma^{-1}\tau)-p},\qquad \sigma,\tau\in S_p,
\end{equation}
where $\cyc$ counts permutation cycles.

At $p=2$, the real formula reduces to
\begin{equation}\label{eq:T2}
 T_2=\begin{pmatrix}
 1&\chi^{-1}&\chi^{-1}\\
 \chi^{-1}&1&\chi^{-1}\\
 \chi^{-1}&\chi^{-1}&1
 \end{pmatrix}.
\end{equation}
This is precisely the fourth-order bond transfer in
\cite[Section~4.4.1, equation~(4.8)]{CEMT}; \eqref{eq:Treal} retains the additional
pairing states required at higher orders.

As matrices, \eqref{eq:Treal}--\eqref{eq:Tcomplex} are $\chi^{-p}$ times the
standard orthogonal and unitary pairing Gram matrices; see
\cite[Section~2]{CollinsMatsumoto}. The real Gram entry is $\chi^{\ell}$,
and restricting to balanced pairings gives the permutation cycle formula.
We use these \emph{Gram matrices}, not their Weingarten pseudoinverses.
No Haar-distributed core or replacement of the Gaussian law is involved.
Their occurrence in the present higher-order rMPS contraction is derived here;
the underlying Gram kernels are not new.

\subsection{A3. Evaluate the row sum without enumerating every path}
The following elementary recurrences evaluate a row sum of the classical
Gram kernels just identified. They are included to make the constants explicit,
not as new combinatorial identities.
\begin{lemma}[Bond-transfer row sum]\label{lem:rowsum}
For every integer $p\ge1$, the transfer matrix in the corresponding field obeys
\begin{equation}\label{eq:rowsum}
 T_p\one=\rho_p(\chi)\one.
\end{equation}
\end{lemma}
\begin{proof}
For a fixed real pairing $\alpha$, define
\[
 F_p(x)=\sum_{\beta\in\mathcal P_2(2p)}x^{\ell(\alpha,\beta)}.
\]
Relabeling positions shows that this does not depend on $\alpha$.
Remove a designated pair $\{u,v\}$ of $\alpha$.
If $\beta$ also pairs $u$ to $v$, the removed component contributes a factor $x$.
Otherwise $\beta$ contains $\{u,a\}$ and $\{v,b\}$.
Delete $u,v$ and replace those two pairs by $\{a,b\}$.
For each smaller pairing there are $p-1$ choices of $\{a,b\}$ and two choices
of its orientation; the number of components is unchanged. Hence
\[
 F_p(x)=(x+2p-2)F_{p-1}(x),\qquad F_1(x)=x.
\]
It follows that
\[
 F_p(\chi)=\chi(\chi+2)\cdots(\chi+2p-2).
\]
Divide by $\chi^p$ to obtain \eqref{eq:rowsum} for real cores.

For complex cores, multiplication by a fixed permutation reduces a row sum to
$\chi^{-p}\sum_{\pi\in S_p}\chi^{\cyc(\pi)}$.
When a new symbol is added to a permutation, it either forms a singleton cycle
or is inserted immediately after one of the $p-1$ existing symbols. Thus the
cycle-count polynomial obeys the recurrence with factor $x+p-1$, giving
\[
 \sum_{\pi\in S_p}x^{\cyc(\pi)}=x(x+1)\cdots(x+p-1).
\]
This proves the complex formula.
\end{proof}

At $p=4$, the two row sums are
\begin{align}
 \rho_4(\chi)&=1+\frac{12}{\chi}+\frac{44}{\chi^2}+\frac{48}{\chi^3}
 =\left(1+\frac2\chi\right)\left(1+\frac4\chi\right)
  \left(1+\frac6\chi\right),&&\F=\R,\label{eq:rho4R}\\
 \rho_4(\chi)&=1+\frac6\chi+\frac{11}{\chi^2}+\frac6{\chi^3}
 =\left(1+\frac1\chi\right)\left(1+\frac2\chi\right)
  \left(1+\frac3\chi\right),&&\F=\C.\label{eq:rho4C}
\end{align}
The real coefficients sum to $105$, and the complex coefficients sum to $24$,
as they should when $\chi=1$.

\subsection{A4. Control arbitrary coefficient tensors and propagate}
\begin{lemma}[Coefficient-network bound]\label{lem:network}
For any fixed choice of the local Wick states, the contraction $D_u$ of the
$2p$ coefficient tensors satisfies
\begin{equation}\label{eq:network}
 |D_u|\le\norm u^{2p}.
\end{equation}
\end{lemma}
\begin{proof}
Each physical index is joined to exactly one other occurrence by a pairing.
The coefficient network therefore has one vertex per copy of $u$ or
$\overline u$, edges joining pairs of vertices, and no self-loops.
For two tensors contracted over all their shared indices, reshape the shared
indices into a single matrix index. Matrix multiplication and Frobenius
Cauchy--Schwarz give
\[
 \norm{\operatorname{contract}(A,B)}_{\mathrm F}
 \le\norm A_{\mathrm F}\norm B_{\mathrm F}.
\]
Merge two vertices, contracting all edges between them, and repeat. At every
step the resulting boundary tensor is retained in full; contractions inside the
merged cluster have already been performed. Each connected component ends as
a scalar bounded by the product of its initial Frobenius norms. Multiplying
these bounds proves \eqref{eq:network}. Complex conjugation does not change the
norm. No assumption about coefficient signs is used.
\end{proof}

Let $c_{\R,p}=(2p-1)!!$ and $c_{\C,p}=p!$.
The Wick expansion in the real case is
\begin{equation}\label{eq:Wick}
 \E|u^*\omega|^{2p}
 =\sum_{\alpha_1,\ldots,\alpha_n}
 \left[\prod_{j=1}^{n-1}T_p(\alpha_j,\alpha_{j+1})\right]
 D_u(\alpha_1,\ldots,\alpha_n).
\end{equation}
The complex expansion has the same structure with permutation labels.
The weights are nonnegative, so Lemmas~\ref{lem:rowsum} and~\ref{lem:network}
give the common scalar endpoint:
\begin{proposition}[Sharp scalar moments]\label{prop:scalar}
For every $u\in\F^{d^n}$ and integer $p\ge1$,
\begin{equation}\label{eq:scalar}
 \boxed{\displaystyle\E|u^*\omega|^{2p}\le c_{\F,p}\mu_p\norm u^{2p}.}
\end{equation}
Equality holds for product tensors.
\end{proposition}
\begin{proof}[Completion of Chain A]
Equation~\eqref{eq:Wick} is at most
\[
 \norm u^{2p}\one^{\mathsf T}T_p^{n-1}\one
 =c_{\F,p}\rho_p(\chi)^{n-1}\norm u^{2p}.
\]
For $u=u_1\otimes\cdots\otimes u_n$, each physical-site contraction is a
product of pairwise inner products of identical vectors (or of conjugate
pairs). Hence every $D_u$ equals $\norm u^{2p}$, and no inequality is lost.
\end{proof}

The $105$ and $24$ local labels at eighth order do not imply that arbitrary
coefficient-network evaluation is computationally cheap. For this uniform
bound, exact evaluation is unnecessary: \eqref{eq:network} handles the
coefficients, and \eqref{eq:rowsum} sums the transfer weights.

\par\Needspace{10\baselineskip}
\section{Chain B: original-law multiplication and conditional energy}\label{sec:B}
\begin{takeaway}
The same moment factor follows from a single Gaussian matrix inequality.
The multiplication and exact-cutoff interfaces locate the relevant original-law
energy; conditional integration evaluates it without enumerating Wick states.
This route can also be written as a standalone conditioning proof.
\end{takeaway}

The fourth-order version of the one-core contraction and its tensor-train
iteration are already in \cite[Section~5.1, Lemmas~3--4]{RR2020}. We generalize
that local estimate in moment order and place the resulting calculation inside
the original-law interfaces of \cite{Framework}. This is not a claim that the
conditional-contraction strategy originated in the framework.

\subsection{B1. Start from the original-law vacuum identity}
Let $\gamma$ be the joint law of the original real Gaussian coordinates in
all cores and probes. In the complex case, use the independent real and
imaginary coordinates with their $1/\sqrt2$ normalization.
On $\mathcal H=\F^2\otimes L^2(\gamma)$, let $\Omega=1$ and let
$\mathcal M_G$ be multiplication by the Hermitian matrix $G$. Then
\begin{equation}\label{eq:vacuum}
 \E\tr G^4=\sum_{b=1}^2\norm{\mathcal M_G^2(e_b\otimes\Omega)}^2.
\end{equation}
This follows directly from $\mathcal M_G^j(e_b\otimes\Omega)=G^je_b$ and
$\sum_b\norm{G^2e_b}^2=\tr G^4$. The dimension factor comes from the
external space $\F^2$, not from the polynomial state space.

\subsection{B2. Use an exact local-degree and active-probe cutoff}
Give the Gaussian polynomial space its Hermite multigrading by the original
cores; the underlying Gaussian chaos and tensor-product structure is classical
\cite[Chapters~2--4]{Janson}. An entry of $y_ay_a^*$ has degree at most two in each core of probe $a$.
After $j\le2$ multiplications, the local degree in every core is at most $2j$,
and a Hermite term has at most $j$ active probe blocks. A block is active when
its local component is nonconstant; this is not the same as having appeared
earlier in an expanded word.

Let $\Pi$ restrict to local core degrees at most four and at most two active
probe blocks. It may also impose the exact invariance under multiplying any
whole core by $-1$, which fixes $G$ and the vacuum. Set
\[
 M=\Pi\mathcal M_G\Pi.
\]
On this finite-dimensional space, multiplication is well defined and
\begin{equation}\label{eq:cutoff}
 M^j(e_b\otimes\Omega)=\mathcal M_G^j(e_b\otimes\Omega),
 \qquad j=0,1,2.
\end{equation}
Thus
\begin{equation}\label{eq:cutoffenergy}
 \E\tr G^4=\sum_{b=1}^2\norm{M^2(e_b\otimes\Omega)}^2.
\end{equation}
This is the exact finite-moment compression interface of v11 \cite{Framework}.
Taking the squared norm involves degree at most eight in each core, which
explains why eighth-order probe information is needed. No large sample outcome
is clipped, and no full-cutoff operator-norm bound is invoked.

\subsection{B3. Retain the coefficient Gram across one core}
In the real case at $p=2$, after transposing the matrix product and matching the
variance normalization, \cite[Lemma~4]{RR2020} gives exactly
\[
 \E_R\norm{RC}_{\mathrm F}^{4}
 =\norm C_{\mathrm F}^{4}+\frac{2}{\chi}\tr\bigl((CC^*)^2\bigr)
 \le\left(1+\frac2\chi\right)\norm C_{\mathrm F}^{4}.
\]
Their proof uses standard Wishart moments. The next proof uses the same
Gaussian Gram geometry and convexity to cover all integer orders and both fields.
\begin{lemma}[One-core Gaussian energy inequality]\label{lem:row}
Let $C\in\F^{s\times t}$ be deterministic, and let
$R\in\F^{\chi\times s}$ have independent Gaussian entries of variance $1/\chi$.
For every integer $p\ge1$,
\begin{equation}\label{eq:row}
 \boxed{\displaystyle\E_R\norm{RC}_{\mathrm F}^{2p}
 \le\rho_p(\chi)\norm C_{\mathrm F}^{2p}.}
\end{equation}
Equality holds when $C$ has rank at most one.
\end{lemma}
\begin{proof}
Let $\lambda_1,\ldots,\lambda_k$ be the nonzero eigenvalues of $CC^*$.
Gaussian orthogonal or unitary invariance gives
\[
 \norm{RC}_{\mathrm F}^2\stackrel d=\sum_{j=1}^k\lambda_jQ_j,
 \qquad Q_j\stackrel d=\frac1\chi\sum_{\ell=1}^{\chi}|g_\ell|^2,
\]
where the $Q_j$ are independent. The gamma density, or its elementary
integration-by-parts recursion, gives
\begin{equation}\label{eq:Qmoments}
 \E Q_j^p=\prod_{l=0}^{p-1}\left(1+\frac{\eta l}{\chi}\right)
 =\rho_p(\chi).
\end{equation}
For real cores, $Q_j$ is $\chi^{-1}$ times a chi-square variable with $\chi$
degrees of freedom. For complex cores it has gamma shape $\chi$ and rate $\chi$.

If $s_0=\sum_j\lambda_j=0$, the assertion is trivial. Otherwise convexity yields
\[
 \left(\sum_j\frac{\lambda_j}{s_0}Q_j\right)^p
 \le\sum_j\frac{\lambda_j}{s_0}Q_j^p.
\]
Take expectations and multiply by $s_0^p$ to obtain \eqref{eq:row}.
With only one nonzero eigenvalue the inequality is an equality.
\end{proof}

Positivity is used only for the eigenvalues of the Gram matrix $CC^*$, not for
the entries of $C$. The estimate therefore applies to arbitrary signed or
complex coefficient tensors. It does not require a positive polynomial cone.

\subsection{B4. Iterate the conditional energy and close the scalar bound}
This is the independent-core/Frobenius-reshaping iteration already used at
fourth order in \cite[Section~5.1]{RR2020}, now with Lemma~\ref{lem:row} at order
$2p$. We give the full argument to specify normalization and equality cases.
\begin{proof}[Independent proof of Proposition~\ref{prop:scalar}]
Start with the coefficient tensor $\overline u$. Reshape its first physical
index as a row index and its remaining physical indices as columns. Its
Frobenius norm is $\norm u$.

Scale the first $n-1$ cores by $\chi^{-1/2}$. Contracting the first core is
exactly the operation $C\mapsto RC$ in Lemma~\ref{lem:row}. Before the next
contraction, combine the current bond index and the next physical index into
the row index. This reshaping preserves the Frobenius norm; all unprocessed
physical indices remain in the column index.

Write $C_j$ for the retained coefficient matrix after the first $j$ cores
and $S_j=\norm{C_j}_{\mathrm F}^2$, with $S_0=\norm u^2$.
The next core is independent of the already processed ones, so
\begin{equation}\label{eq:energyrecursion}
 \E[S_j^p\mid C_{j-1}]\le\rho_p(\chi)S_{j-1}^p,
 \qquad j=1,\ldots,n-1.
\end{equation}
Iterating conditional expectations gives
\[
 \E S_{n-1}^p\le\rho_p(\chi)^{n-1}\norm u^{2p}.
\]
The last unscaled core makes a scalar Gaussian linear form with conditional
variance $S_{n-1}$. Its $2p$-th absolute moment is $c_{\F,p}S_{n-1}^p$.
Taking expectation proves \eqref{eq:scalar}.

For a product tensor, every intermediate coefficient matrix has rank at most
one. All uses of \eqref{eq:row} are then equalities, proving sharpness.
\end{proof}

Each original core is integrated jointly with all its repeated occurrences.
The independence used in \eqref{eq:energyrecursion} is independence of the
\emph{next core} from the preceding ones, not independence of repeated uses
of one core. Only the retained coefficient boundary is reshaped.

The proof of \eqref{eq:scalar} works at all integer orders. For P10 only
$p\le4$ is required. The local-degree cutoff identifies the vacuum energy
needed for the centered matrix calculation, while the conditional contraction
supplies its bond-sensitive moment estimates. The two interfaces are
compatible, but neither their use here nor their notation establishes that
this fixed-order argument requires the specialized graded machinery.

\par\Needspace{10\baselineskip}
\section{Shared completion: projected moments and ordered sample energy}\label{sec:shared}
\begin{takeaway}
The two chains have already converged at Proposition~\ref{prop:scalar}.
The remaining proof is common: convert scalar estimates into projected-vector
moments and expand the centered sample word. Keeping the alternating pairing
is necessary for the exact constants.
\end{takeaway}

\subsection{From scalar projections to the projected norm}
\begin{lemma}[Projected radial moments]\label{lem:radial}
For any deterministic isometry $U:\F^r\to\F^{d^n}$ and integer $p\ge1$,
\begin{equation}\label{eq:radial}
 \E\norm{U^*\omega}^{2p}\le d_{r,p}\mu_p,
 \qquad d_{r,p}=\prod_{j=0}^{p-1}(r+\eta j).
\end{equation}
For $r\le d$, the product-prefix subspace with $r$ last-site directions
attains equality.
\end{lemma}
\begin{proof}
Let $\theta$ be uniform on the unit sphere in $\F^r$. Gaussian polar
decomposition gives, for deterministic $v$,
\[
 \E_\theta|\theta^*v|^{2p}=\frac{c_{\F,p}}{d_{r,p}}\norm v^{2p}.
\]
Indeed, a standard Gaussian vector is its independent radius times its
direction; the radius has $2p$-th moment $d_{r,p}$ and any unit scalar
projection has moment $c_{\F,p}$.
Apply Proposition~\ref{prop:scalar} to $U\theta$, whose norm is one, and use
Fubini. Equality holds on the product-prefix family because every $U\theta$
is a product tensor.
\end{proof}

For $r=2$, writing $R=\norm y^2$, this gives
\begin{equation}\label{eq:radial2}
 \begin{array}{c|cc}
 &\E R^2&\E R^4\\ \hline
 \F=\R&\le8\mu_2&\le384\mu_4\\
 \F=\C&\le6\mu_2&\le120\mu_4
 \end{array}
\end{equation}
The scalar eighth-moment bounds behind the right column are $105\mu_4\norm u^8$
and $24\mu_4\norm u^8$, respectively. These are original eighth moments,
not consequences of matching only four moments.

\subsection{The exact fourth moment of a centered sample average}
The sample-label bookkeeping is the entrywise moment--cumulant principle for
independent sums; see \cite[Chapter~2]{McCullagh}. The following ordered-matrix
specialization follows by direct expansion and is not claimed as a new general
identity. The matrix order must be retained even though each entry is scalar.
\begin{lemma}[Ordered fourth-order sample assembly]\label{lem:assembly}
Let $X_1,\ldots,X_N$ be independent identically distributed centered Hermitian
matrices with finite fourth moments. Write $V=\E X_1^2$ and let $X'$ be an
independent copy. Then
\begin{equation}\label{eq:assembly}
 \E\tr\left(\frac1N\sum_{a=1}^N X_a\right)^4
 =\frac{\E\tr X_1^4+(N-1)
 [2\tr V^2+\E\tr(X_1X'X_1X')]}{N^3}.
\end{equation}
Moreover,
\begin{equation}\label{eq:crossbound}
 \left|\E\tr(X_1X'X_1X')\right|\le\tr V^2.
\end{equation}
\end{lemma}
\begin{proof}
Expand the ordered word of length four in its sample labels. A label appearing
exactly once gives zero after conditioning on the other matrices. There are
$N$ words with one label four times. For each unordered pair of distinct labels,
there are six words: four are cyclically equivalent to $X_1^2(X')^2$, and two
alternate. The former have expected trace $\tr V^2$. Multiplication by
$\binom N2$ and division by $N^4$ give \eqref{eq:assembly}.

For Hermitian $A,B$, Hilbert--Schmidt Cauchy--Schwarz gives
\[
 |\tr(ABAB)|=|\langle BA,AB\rangle_{\mathrm{HS}}|
 \le\norm{BA}_{\mathrm F}\norm{AB}_{\mathrm F}
 =\tr(A^2B^2).
\]
Take expectations and use independence to obtain \eqref{eq:crossbound}.
\end{proof}

The lemma is also the $Q(T)=T^2$ instance of v11's fourth-order original-law
energy identity with component reuse \cite{Framework}. If $K$ is a centered
Gaussian Hermitian matrix with the complete real-entry covariance of $X_1$
(the joint real/imaginary covariance in the complex model), then
\begin{equation}\label{eq:calibration}
 \E\tr G^4=\frac{N-1}{N^3}\E\tr K^4+\frac1{N^3}\E\tr X_1^4.
\end{equation}
Here $\E\tr K^4=2\tr V^2+\E\tr(X_1X'X_1X')$.
This is a comparison of the centered matrices $X_a$, not a substitution of
Gaussian surrogates for the original tensor cores.

\subsection{Proof of the uniform endpoint}
\begin{proof}[Proof of Theorem~\ref{thm:uniform}]
For $Y=yy^*$, $X=Y-I_2$, and $R=\norm y^2$, the eigenvalues of $X$ are
$R-1$ and $-1$. Since isotropy gives $\E R=2$,
\[
 \E\tr X^4=\E R^4-4\E R^3+6\E R^2-6.
\]
With an independent copy $R'$, the inequality
$\E[(R^2-(R')^2)(R-R')]\ge0$ implies
$\E R^3\ge(\E R^2)(\E R)=2\E R^2$. Hence
\begin{equation}\label{eq:Aupper}
 \E\tr X^4\le\E R^4.
\end{equation}
Also $V=\E X^2\succeq0$, so
\begin{equation}\label{eq:Vupper}
 \tr V^2\le(\tr V)^2=(\E R^2-2)^2.
\end{equation}
Substitute \eqref{eq:radial2}, \eqref{eq:crossbound}, \eqref{eq:Aupper}, and
\eqref{eq:Vupper} into \eqref{eq:assembly}. This gives
\eqref{eq:uniformR} and \eqref{eq:uniformC} with exactly the constants stated.
\end{proof}

Neither scalar route requires independence of the coordinates of $y$. The
sample-label identity uses only independence of different probes, while their
within-probe correlations are retained in the moments of $X$.

\par\Needspace{10\baselineskip}
\section{The common exact product-prefix calculation}\label{sec:exact}
\begin{takeaway}
For the specified subspace $U_0$, the projected vector has the exact original-law
representation $y\stackrel d=\sqrt W\,z$. Both coordinates share $W$.
This closes the radial moment boundary and permits exact evaluation of the
alternating sample word, not only a bound on it.
\end{takeaway}

\subsection{A closed radial boundary from the original cores}
For $U_0$ in \eqref{eq:prefixU}, define the prefix row states
\[
 v_1=\chi^{-1/2}H_1(1,:),\qquad
 v_j=\chi^{-1/2}v_{j-1}H_j(:,1,:),\quad 2\le j\le n-1.
\]
The final output is
\[
 y=\bigl(v_{n-1}H_n(:,1),\;v_{n-1}H_n(:,2)\bigr)^{\mathsf T}.
\]
Conditionally on preceding cores, the coordinates of $v_j$ are independent
Gaussians with variance $\norm{v_{j-1}}^2/\chi$. Thus the ratio
$\norm{v_j}^2/\norm{v_{j-1}}^2$ has the law of $Q$ in \eqref{eq:Qmoments},
independent of the preceding sigma-field. The denominator is nonzero almost
surely. The first factor has the same law, and the last core supplies a
normalized Gaussian vector independent of these ratios. Consequently
\begin{equation}\label{eq:mixture}
 \boxed{\displaystyle
 y\stackrel d=\sqrt W\,z,\qquad
 W=\prod_{j=1}^{n-1}Q_j,\qquad z\sim N_{\F}(0,I_2),}
\end{equation}
where $Q_1,\ldots,Q_{n-1},z$ are independent and $\E W^p=\mu_p$.
This is derived from the full core law. It is not a decoupling assumption.

On moment functions through order four, the radial conditional-integration
operator has the exact update
\begin{equation}\label{eq:radialtransfer}
 (\mathsf T_\chi f)(s)=\E f(sQ),\qquad
 \mathsf T_\chi(s^p)=\rho_p(\chi)s^p.
\end{equation}
Its matrix in the basis $1,s,s^2,s^3,s^4$ is
$\diag(1,1,\rho_2,\rho_3,\rho_4)$.
This is a five-dimensional closure for these moment functions on the prefix
family, not for the full probability law or arbitrary tensor coefficients.
The Wick route obtains the same moments because every coefficient contraction
for a product direction attains equality.

\subsection{The one-probe fourth moment}
Let $s=\norm z^2$, so $R=Ws$ and
\[
 \E s^p=\prod_{j=0}^{p-1}(2+\eta j).
\]
For $a=\mu_2$, $b=\mu_3$, $c=\mu_4$, direct expansion gives
\begin{align}
 A_{\F}:=\E\tr X^4
 &=\E R^4-4\E R^3+6\E R^2-6,\label{eq:Aexact}\\
 V:=\E X^2&=vI_2,\qquad v=(2+\eta)a-1.\label{eq:Vexact}
\end{align}
For the real field, $(\E s^2,\E s^3,\E s^4)=(8,48,384)$; for the complex
field they are $(6,24,120)$. Substitution proves the values of $A_{\F}$ in
\eqref{eq:ABR}--\eqref{eq:ABC}, and $\E\tr G^2=2v/N$ proves \eqref{eq:second}.

\subsection{The alternating word and the exact pair coefficient}
For deterministic real symmetric or complex Hermitian $D$, Gaussian
fourth-moment contraction yields
\[
 \E[zz^*Dzz^*]=\eta D+(\tr D)I_2.
\]
Expanding $XDX=(Wzz^*-I_2)D(Wzz^*-I_2)$ therefore gives
\begin{equation}\label{eq:channel}
 \mathcal L(D):=\E[XDX]=(\eta a-1)D+a(\tr D)I_2.
\end{equation}
In addition,
\[
 \E\tr X^2=2v,\qquad \E(\tr X)^2=2(2+\eta)a-4.
\]
Condition on $X$ and integrate $X'$ to obtain
\begin{align}
 C_{\F}:=\E\tr(XX'XX')
 &=\E\tr\bigl(X\mathcal L(X)\bigr)\notag\\
 &=(\eta a-1)\,2v+a\bigl(2(2+\eta)a-4\bigr).\label{eq:Cexact}
\end{align}
Explicitly,
\[
 C_{\R}=24a^2-16a+2,\qquad C_{\C}=12a^2-12a+2.
\]
Since $2\tr V^2=4v^2$, the pair coefficient is
\[
 B_{\F}=4v^2+C_{\F}.
\]
This gives $B_{\R}=88a^2-48a+6$ and $B_{\C}=48a^2-36a+6$.
Inserting these and \eqref{eq:Aexact} into Lemma~\ref{lem:assembly} completes
the proof of Theorem~\ref{thm:exact}.

\subsection{A normalization check}
As $\chi\to\infty$ with $n$ fixed, $\mu_p\to1$. Both chains therefore give
\begin{equation}\label{eq:Gaussianlimit}
 \E\tr G^4\longrightarrow
 \begin{cases}
 46/N^2+188/N^3,&\F=\R,\\
 18/N^2+36/N^3,&\F=\C.
 \end{cases}
\end{equation}
The one-probe and pair constants in this limit are $(234,46)$ and $(54,18)$,
respectively. These follow by direct substitution, without a limiting
interchange of unbounded random variables.

\par\Needspace{24\baselineskip}
\section{End results and explicit constants}\label{sec:comparison}
\subsection{The same constants, not merely the same order}
\begingroup\small\setlength{\tabcolsep}{5pt}\renewcommand{\arraystretch}{1.25}
\begin{tabularx}{\textwidth}{@{}P{.33\textwidth}YY@{}}
\toprule
\textbf{Quantity}&\textbf{Chain A: Wick / transfer}&\textbf{Chain B: multiplication / energy}\\
\midrule
One-bond moment factor&$\rho_p(\chi)$ from the exact row sum&$\rho_p(\chi)$ from conditional energy\\
Length multiplier&$\mu_p=\rho_p(\chi)^{n-1}$&$\mu_p=\rho_p(\chi)^{n-1}$\\
Real scalar eighth-moment allowance for $\norm u=1$&$105\mu_4$&$105\mu_4$\\
Complex scalar eighth-moment allowance for $\norm u=1$&$24\mu_4$&$24\mu_4$\\
Projected real eighth-moment allowance, dimension two&$384\mu_4$&$384\mu_4$\\
Projected complex eighth-moment allowance, dimension two&$120\mu_4$&$120\mu_4$\\
Uniform Gram fourth moment&\eqref{eq:uniformR} and \eqref{eq:uniformC}&Exactly the same bounds\\
Product-prefix Gram fourth moment&\eqref{eq:exact}--\eqref{eq:ABC}&Exactly the same formula\\
\bottomrule
\end{tabularx}
\endgroup

The scalar and projected radial bounds are sharp because product directions
or the product-prefix subspace attain them. The \emph{uniform centered Gram}
bound need not be sharp: its common completion uses
$\tr V^2\le(\tr V)^2$, bounds the alternating word in absolute value, and
replaces $\E\tr X^4$ by $\E\norm y^8$. These losses are shared by both chains.


\subsection{Why the constants coincide: the same rank-one Gaussian extremum}
\label{subsec:equalconstants}
Equal constants do not follow merely from studying the same random law: two
estimates of that law could lose different amounts. Here both routes preserve
the exact same local Gaussian moment and then take the same worst-case
coefficient envelope. The extremizer belongs to the class under consideration.

\step{The Wick row sum is a radial Gaussian moment.}
Let $g\in\F^\chi$ be standard Gaussian and put $Q=\norm g^2/\chi$.
Expanding $(\sum_i g_i^2)^p$ in the real case fixes one pairing $\alpha$ that
identifies the two appearances of each squared coordinate. Applying Isserlis'
formula adds the second pairing $\beta$. Every component of
$\alpha\cup\beta$ contributes one free coordinate index. Thus
\[
 \underbrace{\sum_{\beta\in\mathcal P_2(2p)}
                  \chi^{\ell(\alpha,\beta)-p}}_{\text{Chain A: one row sum}}
 =\underbrace{\E\left(\frac{\norm g^2}{\chi}\right)^p}_{\text{Chain B: gamma moment}}
 =\rho_p(\chi).
\]
The complex identity uses balanced pairings, equivalently permutations, in
place of $\mathcal P_2(2p)$. It has the same interpretation. This is the exact
link between the two calculations, not an analogy between their notation.
The Wick identities \cite{Isserlis,Goodman} and the pairing Gram representation
\cite{CollinsMatsumoto} give the two ways of reading this equality.

\step{Both coefficient bounds are saturated by rank one.}
For $K=CC^*\succeq0$, let $s=\tr K$ and $S=\norm{RC}_{\mathrm F}^2$.
If $\lambda_j$ are the eigenvalues of $K$, then
$S\stackrel d=\sum_j\lambda_jQ_j$. The convexity step in Chain~B gives
\[
 \sup_{C:\,\norm C_{\mathrm F}^2=s}\E S^p
       =s^p\rho_p(\chi),
\]
where any admissible rank-one $C$ attains the supremum. For a product coefficient
tensor, every intermediate matrix has rank one, while every coefficient
network in Chain~A attains its Frobenius bound. Equality therefore persists
through all $n-1$ bonds in both proofs. The scalar factor
$c_{\F,p}\mu_p$ cannot be lowered in a theorem uniform over all $u$.

\step{A trace-power calculation exposes the common information loss.}
This also follows from the classical trace--cumulant description of Gaussian
quadratic forms \cite{Lancaster}; the complex setting is treated in
\cite{Goodman}. Write $\tau_k=\tr K^k$, so $\tau_1=s$, and set
$\gamma_\chi=\eta/\chi$. Diagonalization and the gamma moment-generating
function give, for $t$ near zero,
\[
 \E e^{tS}=\det(I-\gamma_\chi tK)^{-1/\gamma_\chi},
 \qquad
 \kappa_k(S)=(k-1)!\gamma_\chi^{k-1}\tau_k.
\]
The second formula is obtained by expanding the logarithm of the first; no new
quadratic-form cumulant theorem is needed. Substituting these cumulants into
the fourth moment--cumulant identity \cite[Chapter~2]{McCullagh} yields
\[
 \begin{aligned}
 \E S^4={}&s^4+6\gamma_\chi s^2\tau_2
     +3\gamma_\chi^2\tau_2^2
     +8\gamma_\chi^2s\tau_3+6\gamma_\chi^3\tau_4.
 \end{aligned}
\]
All coefficients are nonnegative. Replacing $\tau_k$ by $s^k$ gives
\[
 \E S^4\le s^4(1+6\gamma_\chi+11\gamma_\chi^2+6\gamma_\chi^3)
       =s^4\rho_4(\chi).
\]
At rank one, $\tau_k=s^k$, so the bound is exact. With $\eta=2$ this is the
Wick row polynomial $(1,12,44,48)$, and with $\eta=1$ it is
$(1,6,11,6)$; compare \eqref{eq:rho4R}--\eqref{eq:rho4C}.
The loop expansion and the spectral-energy calculation thus collapse the same
positive Gaussian contraction to the same rank-one envelope.

\step{The equality continues because the rest of the proof is shared.}
Spherical averaging contributes the same exact factor $d_{r,p}/c_{\F,p}$.
The ordered sample-label identity is exact, and the subsequent three upper
bounds described above are the same in both chains. The equality of the
constants tables follows from this common completion, not from a theorem that
all Wick and graded proofs must have equal constants.

There is also an operator interpretation. The symmetric nonnegative transfer
has $T_p\one=\rho_p\one$ and operator norm $\rho_p$. On the rank-one boundary,
conditional integration sends $s^p$ to $\E(sQ)^p=\rho_ps^p$. These are two
representations of the same invariant scalar moment. For arbitrary coefficient
matrices, the radial description is an upper envelope, not an exact compression
of their full law or an equivalence of the full operators.

\step{What would be needed to improve a bound.}
With a spread-out coefficient Gram, $\tr K^2<s^2$ and some higher trace powers
are also strictly smaller than their rank-one values. Retaining these invariants
can improve a coefficient-sensitive estimate in either language. A useful
multi-site refinement must also justify how those invariants are updated after
the next core contraction and reshaping. The worst-case scalar theorem cannot
improve uniformly, because product tensors attain it; the uniform centered-Gram
bound may still improve by avoiding its separate shared relaxations.

\subsection{Tensor length versus bond dimension}
Put $\tau=(n-1)/\chi$. The inequality $\log(1+x)\le x$ gives
\begin{equation}\label{eq:muexp}
 \mu_p\le\exp\left(\frac{\eta p(p-1)}2\tau\right).
\end{equation}
Consequently, both uniform theorems imply
\begin{align}
 \E\tr G^4&\le\frac{192e^{4\tau}}{N^2}
                    +\frac{384e^{12\tau}}{N^3},&&\F=\R,\label{eq:uniformexpR}\\
 \E\tr G^4&\le\frac{108e^{2\tau}}{N^2}
                    +\frac{120e^{6\tau}}{N^3},&&\F=\C.\label{eq:uniformexpC}
\end{align}
For the product-prefix family the leading coefficients improve:
\begin{align}
 \E\tr G^4&\le\frac{88e^{4\tau}}{N^2}
                    +\frac{384e^{12\tau}}{N^3},&&\F=\R,\\
 \E\tr G^4&\le\frac{48e^{2\tau}}{N^2}
                    +\frac{120e^{6\tau}}{N^3},&&\F=\C.
\end{align}
Here $\mu_3\ge\mu_2\ge1$ implies $A_{\R}\le384\mu_4$ and
$A_{\C}\le120\mu_4$, while $B_{\R}\le88\mu_2^2$ and
$B_{\C}\le48\mu_2^2$. If $\chi$ is proportional to $n$, all these fixed-order
constants remain bounded uniformly in the tensor length.

For comparison, a total-degree Gaussian hypercontractive estimate for the real
scalar polynomial $u^*\omega$ gives
\[
 \norm{u^*\omega}_{L_8}\le7^{n/2}\norm u,
 \qquad \E|u^*\omega|^8\le7^{4n}\norm u^8.
\]
This is the degree-$n$ consequence of Nelson's Gaussian hypercontractivity
\cite{Nelson}; see also \cite[Chapter~5]{Janson} and the later reproof
\cite[equation~(1.8) of the arXiv version]{DPLS}. The latter is a source for a
proof, not the origin of the inequality. Both chains improve the bond-blind
allowance to $105\mu_4\le105e^{12\tau}$, so this gain is not a quantitative
separation between Chains A and B.

\subsection{The probability certificates being compared}
The spectral theorem and Markov's inequality give
\begin{equation}\label{eq:Markov}
 \Prob\{\norm G>\varepsilon\}\le\varepsilon^{-4}\E\tr G^4.
\end{equation}
For the prefix family, define
\[
 (A,B)=(A_{\F},B_{\F});
\]
for the uniform bound, define
\begin{equation}\label{eq:uniformAB}
 (A,B)=\begin{cases}
 (384\mu_4,\;3(8\mu_2-2)^2),&\F=\R,\\
 (120\mu_4,\;3(6\mu_2-2)^2),&\F=\C.
 \end{cases}
\end{equation}
In either case, the explicit sufficient certificate is
\begin{equation}\label{eq:certificate}
 A+(N-1)B\le\delta\varepsilon^4N^3.
\end{equation}
A simpler, nonminimal sufficient condition is
\begin{equation}\label{eq:simplercert}
 N\ge\max\left\{\left(\frac{2A}{\delta\varepsilon^4}\right)^{1/3},
                      \left(\frac{2B}{\delta\varepsilon^4}\right)^{1/2}\right\}.
\end{equation}
This follows by bounding the fourth moment by $A/N^3+B/N^2$.
For fixed $\varepsilon$ and $\tau$, its leading high-confidence dependence is
$\delta^{-1/2}$. This is not logarithmic confidence dependence or a claim of
optimal sample complexity.

The second-moment prefix certificate used as a baseline is
\begin{equation}\label{eq:secondcert}
 N\ge\frac{h_{\F}\mu_2-2}{\delta\varepsilon^2},
 \qquad h_{\R}=8,\quad h_{\C}=6.
\end{equation}
It concerns the \emph{same} subspace $U_0$. The comparison is between these
specified moment certificates, not all algorithms or all concentration methods
available for rMPS.

\subsection{Concrete sufficient probe counts}
All entries in Table~\ref{tab:counts} use $n=100$, $\varepsilon=1/2$, and
$\delta=0.01$. Each is the smallest integer satisfying the corresponding
explicit certificate. The exact fourth-moment columns are repeated to make the
absence of a difference between proof chains visible.

\begin{table}[htbp]\centering
\small\setlength{\tabcolsep}{5pt}\renewcommand{\arraystretch}{1.25}
\begin{tabularx}{\textwidth}{@{}P{.13\textwidth}r
 >{\raggedleft\arraybackslash}X >{\raggedleft\arraybackslash}X
 >{\raggedleft\arraybackslash}X >{\raggedleft\arraybackslash}X@{}}
\toprule
& &\textbf{Second moment}&\multicolumn{2}{c}{\textbf{Exact fourth moment}}&\textbf{Uniform fourth}\\
\cmidrule(lr){4-5}
\textbf{Field}&\textbf{$\chi$}&\textbf{Prefix $U_0$}&\textbf{Chain A}&\textbf{Chain B}&\textbf{Both chains}\\
\midrule
Real&100&21\,929&4\,611&4\,611&5\,234\\
Real&400&4\,444&531&531&779\\
Real&1\,000&3\,100&360&360&540\\
Complex&100&5\,628&706&706&1\,009\\
Complex&400&2\,274&253&253&396\\
Complex&1\,000&1\,850&201&201&322\\
\bottomrule
\end{tabularx}
\caption{Sufficient probe counts from \eqref{eq:secondcert} and
\eqref{eq:certificate}. The uniform column is valid for every fixed
two-dimensional $U$. These are theorem certificates, not empirical thresholds.}
\label{tab:counts}
\end{table}

Exact rational arithmetic checks both the passing integer and the preceding
failing integer; Appendix~\ref{app:checks} records the procedure.
For real cores with $n=100$, $\chi=400$, $\varepsilon=1/2$, and
$\delta=10^{-6}$, the same calculations give
\[
 \begin{array}{c|r}
 \text{Certificate}&\text{Sufficient }N\\ \hline
 \text{Second moment, }U_0&44\,431\,236\\
 \text{Exact fourth moment, }U_0\text{ (either chain)}&51\,182\\
 \text{Uniform fourth moment (either chain)}&76\,967
 \end{array}
\]
The improvement here is from using the fourth moment rather than the second,
not from choosing Chain B instead of Chain A.

\begin{takeaway}
\textbf{Comparison outcome.} At every displayed parameter choice, the two
fourth-order proof chains yield the same sufficient probe counts because they
produce the same inequalities. Their difference is the retained description and
proof organization, not a smaller numerical constant.
\end{takeaway}

\par\Needspace{10\baselineskip}
\section{Scope, provenance, and what does not extend unchanged}\label{sec:scope}
\subsection{A fourth-moment surrogate is not an eighth-moment identity}
At fourth order, the real transfer has the special form
\[
 T_2=\left(1-\frac1\chi\right)I+\frac1\chi\one\one^{\mathsf T}.
\]
This underlies the source paper's fractured-Gaussian fourth-moment surrogate
\cite[Section~4.4.2, Theorem~4.4]{CEMT}. At eighth order, the transfer has distinct off-diagonal weights
$\chi^{-1}$, $\chi^{-2}$, and $\chi^{-3}$, and for $\chi>1$ it is not the same
two-term matrix form.

On a unit real product coordinate, the surrogate's segmentation gives a
product of $1+\operatorname{Binomial}(n-1,1/\chi)$ standard Gaussians.
Its eighth moment is therefore
\begin{equation}\label{eq:surrogate8}
 105\left(1+\frac{104}{\chi}\right)^{n-1}.
\end{equation}
The actual rMPS eighth moment, computed by either chain, is
\begin{equation}\label{eq:original8}
 105\left[\left(1+\frac2\chi\right)
           \left(1+\frac4\chi\right)
           \left(1+\frac6\chi\right)\right]^{n-1}.
\end{equation}
These generally differ. The precise limitation is that the existing
\emph{surrogate} cannot be reused unchanged. The original-core Wick method
itself extends successfully, as Section~\ref{sec:A} shows.

\subsection{What the framework contributes in this comparison}
Chain B uses the original-law multiplication identity, a moment-preserving
reachable-state cutoff, a valid coefficient-boundary energy estimate, and the
fourth-order sample energy assembly. It does not use the positive-coefficient
cone theorem, connected-creation graph estimates, shared-factor occupation
bounds, or covariance-spectrum renewal.

The one-core inequality can be proved independently by rotational invariance
and convexity; its real fourth-order instance and tensor-train iteration have
the direct published precedent \cite[Section~5.1]{RR2020}. The centered-sample
identity follows independently by expanding four labels, or by the classical
moment--cumulant calculus \cite{McCullagh}. The framework supplies a consistent
organization and precise original-law interfaces, without presenting these
Gaussian and cumulant identities as new. This pilot does not establish that its specialized
operator machinery is indispensable or stronger than the extended Wick method.

The distinction is substantive: Chain~B generalizes and reorganizes an existing
Gaussian contraction strategy, and Chain~A extends an existing Wick/transfer
strategy. Section~\ref{subsec:equalconstants} identifies the shared rank-one
extremum that explains their identical numerical constants.

\subsection{The boundary of the result}
The fixed-order P10 calculation supplies a uniform bound for arbitrary fixed
two-dimensional subspaces and an exact formula on the specified coupled family.
It does not supply a sharp theorem for growing subspace dimension, a prescribed
optimal bond-dimension tradeoff, or logarithmic dependence on failure probability.
Although the scalar bound holds at every integer order, the centered matrix
assembly here has only been carried out at order four.

The note consolidates the rMPS proof and the subsequent two-method comparison
from this project. It is an analytic derivation, not an independently reviewed
proof, a formal proof-assistant certificate, or a publication-priority claim.
Gaussian tensor-train projection and fourth-order conditional-contraction
antecedents are credited to \cite{RR2020}; the exact fourth-moment rMPS transfer
and surrogate analysis are credited to \cite{CEMT}. Standard pairing Gram
matrices, Gaussian moment identities, and cumulant bookkeeping are credited at
their points of use. The multiplication and energy interfaces are credited to
\cite{Framework}, whose bibliographic author and title follow the supplied file.
The higher-order comparison and P10 constants are proved here without asserting
that no equivalent result exists elsewhere.

\par\Needspace{12\baselineskip}
\appendix
\section{Certificate checks and reproducibility}\label{app:checks}
\subsection{Exact arithmetic for the reported probe counts}
For integer $n,\chi$ and $p\le4$, compute
\[
 \mu_p=\left[\prod_{j=0}^{p-1}\frac{\chi+\eta j}{\chi}\right]^{n-1}
\]
as a rational number, not from a floating-point exponential approximation.
Use these exact values to form $A,B$ from \eqref{eq:ABR}, \eqref{eq:ABC}, or
\eqref{eq:uniformAB}. For rational $\varepsilon,\delta$, define
\[
 F(N)=\frac{A+(N-1)B}{N^3},\qquad H=\delta\varepsilon^4.
\]
The table is generated by doubling an upper bracket until $F(N)\le H$, then
binary-searching the bracket. Both $F(N)\le H$ and $F(N-1)>H$ are checked
exactly at the returned integer.

For the listed parameter sets, $3A\ge B\ge0$ is also checked exactly. This
justifies monotonicity on $N\ge1$, since
\[
 F'(x)=\frac{-2Bx-3(A-B)}{x^4}\le0\qquad(x\ge1).
\]
The second-moment entry is the rational ceiling in \eqref{eq:secondcert}.
Thus these entries are minima of the \emph{specified certificates}, not minima
of the true probability of embedding failure.

\subsection{Independent normalization and original-core tests}
The previously supplied \texttt{rmps\_p10\_checks.py} implements these rational
calculations and tensor Gauss--Hermite consistency tests for small original-core
instances. It includes real and complex cores, lengths two and three, bond
dimension two, and a non-product real two-dimensional coefficient subspace.
The isotropy, one-probe fourth coefficient, and alternating-pair coefficient
are checked separately.

The degree-$2k-1$ exactness of $k$-node Gauss--Hermite quadrature is classical;
see Golub--Welsch \cite{GolubWelsch} for its Jacobi-matrix construction.
Five nodes integrate a real standard Gaussian polynomial through degree nine,
whereas these test integrands have degree at most eight in each primitive
coordinate. The implementation evaluates quadrature in floating point,
so those tests are numerical consistency checks, not exact-arithmetic proof
certificates. The rational probe-count checks are exact. The analytic proofs
above do not depend on either finite test family.

Additional transfer checks enumerate the $105$ real pairings and $24$ complex
permutations at eighth order. They give the component-count vectors
$(1,12,44,48)$ and $(1,6,11,6)$, respectively, agreeing with
\eqref{eq:rho4R}--\eqref{eq:rho4C}. The Gaussian-limit constants
$(A_{\R},B_{\R})=(234,46)$ and $(A_{\C},B_{\C})=(54,18)$ are verified by direct
substitution. These finite checks validate normalization and implementation,
not all parameter values of the theorem.

\clearpage
\begin{thebibliography}{99}
\raggedright
\addcontentsline{toc}{section}{References}
\setlength{\itemsep}{8pt}

\bibitem{CEMT}
C.~Cama\~no, E.~N.~Epperly, R.~A.~Meyer, and J.~A.~Tropp,
\emph{Linear algebra at exponential scale via tensor network dimension reduction},
2026. \href{https://arxiv.org/abs/2606.15350v1}{arXiv:2606.15350v1}.
\textit{Role:} the model normalization in Definition~1.1; the fourth-order
transfer in Section~4.4.1, equation~(4.8), and Theorem~4.3; the fourth-moment
fractured-Gaussian surrogate in Section~4.4.2, Theorem~4.4; the complex version
in Section~4.6. The higher-order calculation here is an extension, not a theorem
attributed to that paper.

\bibitem{RR2020}
B.~T.~Rakhshan and G.~Rabusseau,
\emph{Tensorized Random Projections},
Proceedings of the Twenty Third International Conference on Artificial
Intelligence and Statistics, Proceedings of Machine Learning Research
\textbf{108} (2020), 3306--3316.
\href{https://proceedings.mlr.press/v108/rakhshan20a.html}{Publisher record and paper}.
\textit{Role:} Gaussian tensor-train projection precedent; Theorem~1 and
Section~5.1, especially Lemmas~3--4, give the direct fourth-order predecessor of
our one-core contraction and independent-core iteration. The all-order
extension is proved in this note.

\bibitem{Isserlis}
L.~Isserlis,
\emph{On a formula for the product-moment coefficient of any order of a normal
frequency distribution in any number of variables},
Biometrika \textbf{12}(1--2) (1918), 134--139.
\href{https://doi.org/10.1093/biomet/12.1-2.134}{doi:10.1093/biomet/12.1-2.134}.
\textit{Role:} the classical probabilistic Gaussian pairing identity, often
called Wick's rule. It is the local Gaussian input to Chain~A.

\bibitem{Goodman}
N.~R.~Goodman,
\emph{Statistical analysis based on a certain multivariate complex Gaussian
distribution (an introduction)},
The Annals of Mathematical Statistics \textbf{34}(1) (1963), 152--177.
\href{https://doi.org/10.1214/aoms/1177704250}{doi:10.1214/aoms/1177704250}.
\textit{Role:} the classical complex Gaussian/Wishart setting, including the
circular covariance structure. Balanced pairings here also follow directly
from Isserlis applied to real and imaginary parts.

\bibitem{CollinsMatsumoto}
B.~Collins and S.~Matsumoto,
\emph{On some properties of orthogonal Weingarten functions},
Journal of Mathematical Physics \textbf{50}(11) (2009), 113516.
\href{https://doi.org/10.1063/1.3251304}{doi:10.1063/1.3251304};
\href{https://arxiv.org/abs/0903.5143}{arXiv:0903.5143}.
\textit{Role:} Section~2 records the orthogonal pairing Gram kernel
$d^{\mathrm{loop}}$ and its unitary balanced-pairing counterpart. Our transfer
matrices are normalized versions of these classical Grams; no Weingarten
pseudoinverse or Haar integration is used in this proof.

\bibitem{Lancaster}
H.~O.~Lancaster,
\emph{Traces and cumulants of quadratic forms in normal variables},
Journal of the Royal Statistical Society, Series~B (Methodological)
\textbf{16}(2) (1954), 247--254.
\href{https://doi.org/10.1111/j.2517-6161.1954.tb00167.x}{doi:10.1111/j.2517-6161.1954.tb00167.x}.
\textit{Role:} the trace-power/cumulant description of real Gaussian quadratic
forms, used to explain the equality of the constants. The determinant formula
and the complex specialization are derived explicitly in this note.

\bibitem{McCullagh}
P.~McCullagh,
\emph{Tensor Methods in Statistics}, Chapman and Hall, London, 1987.
\href{https://www.stat.uchicago.edu/~pmcc/tensorbook/}{Author-hosted first-edition chapters}.
\textit{Role:} Chapter~2, the elementary moment--cumulant calculus underlying
independent-label assembly. Our ordered-matrix identity is proved entrywise
and is not claimed as a new general cumulant result.

\bibitem{Janson}
S.~Janson,
\emph{Gaussian Hilbert Spaces}, Cambridge Tracts in Mathematics, vol.~129,
Cambridge University Press, 1997.
\href{https://doi.org/10.1017/CBO9780511526169}{doi:10.1017/CBO9780511526169}.
\textit{Role:} classical Gaussian chaos, Wick/tensor structure, and
hypercontractivity; see Chapters~2--5. This is background for the Hermite
grading, not attribution of a new cutoff principle.

\bibitem{Nelson}
E.~Nelson,
\emph{The free Markoff field},
Journal of Functional Analysis \textbf{12}(2) (1973), 211--227.
\href{https://doi.org/10.1016/0022-1236(73)90025-6}{doi:10.1016/0022-1236(73)90025-6}.
\textit{Role:} classical Gaussian hypercontractivity, used only for the
bond-blind total-degree comparison, not for the main rMPS contraction.

\bibitem{DPLS}
P.~Da Pelo, A.~Lanconelli, and A.~I.~Stan,
\emph{H\"older--Young--Lieb inequalities for norms of Gaussian Wick products},
Infinite Dimensional Analysis, Quantum Probability and Related Topics
\textbf{14}(3) (2011), 375--407.
\href{https://doi.org/10.1142/S0219025711004456}{doi:10.1142/S0219025711004456};
\href{https://arxiv.org/abs/1101.2947}{arXiv:1101.2947}.
\textit{Role:} a later proof of Nelson hypercontractivity; equation~(1.8) in
the arXiv version. The published title is plural; the preprint is titled
\emph{A H\"older--Young--Lieb inequality for norms of Gaussian Wick products}.
Neither version is credited with originating hypercontractivity.

\bibitem{GolubWelsch}
G.~H.~Golub and J.~H.~Welsch,
\emph{Calculation of Gauss quadrature rules},
Mathematics of Computation \textbf{23}(106) (1969), 221--230.
\href{https://doi.org/10.1090/S0025-5718-69-99647-1}{doi:10.1090/S0025-5718-69-99647-1}.
\textit{Role:} the classical Jacobi-matrix construction of Gaussian quadrature
rules, relevant only to the numerical consistency checks.

\bibitem{Framework}
D.~Heidary,
\emph{Graded Multiplication Operators for Random-Matrix Moments},
Version~11: exposition and attribution revision, 16 September 2026.
Supplied unpublished manuscript \texttt{v11\_attribution\_revised.tex}.
\textit{Role:} exact finite moment compression and the original-law
quadratic-energy interface with component reuse. The author and title are
transcribed from the supplied manuscript. These organizing interfaces do not
replace the public Gaussian contraction, pairing, or cumulant precedents above.
\end{thebibliography}

\end{document}
